\chapter{Identitäten und Gleichungslogik}

\section{Wortidentitäten und HSP-Konstruktion}

\begin{definition}[Wortidentität und Gleichungsklasse]
Es sei $X=\{x_1,x_2,\ldots\}$ eine feste abzählbar unendliche
Variablenmenge. Eine \emph{Wortidentität} ist ein Paar $(u,v)\in X^+\times X^+$,
geschrieben $u=v$. Eine Halbgruppe $S$ erfüllt sie, wenn jede Belegung
$X\to S$ nach Auswertung der Wörter denselben Wert für $u$ und $v$ liefert.
\emph{Identität} bedeutet hier Wortgleichung. Die Wörter sind nichtleer;
als Operation wird nur die Multiplikation verwendet, ohne
Einheitssymbol oder Inversenoperation.

Für eine Klasse $\mathcal K$ von Halbgruppen bezeichnet
$\operatorname{Id}(\mathcal K)$ die Menge ihrer gemeinsamen
Wortidentitäten. Für eine Menge $\Sigma$ von Wortidentitäten ist
$\operatorname{Mod}(\Sigma)$ die Klasse aller Halbgruppen, die sämtliche
Identitäten aus $\Sigma$ erfüllen. Eine solche Klasse heißt
\emph{Gleichungsklasse}.
\end{definition}

\begin{definition}[HSP-Konstruktion]
$\operatorname{HSP}(\mathcal K)$ ist die Klasse aller Halbgruppen,
die bis auf Isomorphie homomorphe Bilder von Unterhalbgruppen
mengenindizierter direkter Produkte von Halbgruppen aus $\mathcal K$ sind.
Leere Unterhalbgruppen sind zugelassen;
das Produkt der leeren Familie ist die einelementige Halbgruppe.
\end{definition}

\begin{theorem}[Birkhoffs HSP-Satz für Halbgruppen]
\label{b28guide:thm:hsp}
Für jede Klasse $\mathcal K$ von Halbgruppen gilt
\[
 \operatorname{Mod}\bigl(\operatorname{Id}(\mathcal K)\bigr)
 =\operatorname{HSP}(\mathcal K).
\]
Dies ist zugleich die kleinste Klasse, die $\mathcal K$ enthält und
unter Isomorphie, homomorphen Bildern, Unterhalbgruppen und beliebigen
mengenindizierten Produkten einschließlich des leeren Produkts
abgeschlossen ist.

Eine Klasse von Halbgruppen ist genau dann eine Gleichungsklasse,
wenn sie diese Abschlußeigenschaften besitzt.
Für $\mathcal K=\varnothing$ bestehen beide Seiten der Gleichung
aus den Halbgruppen mit höchstens einem Element.
\cite{birkhoff1935,burris}
\end{theorem}

\section{Vollständigkeit der Gleichungslogik}

\begin{definition}[Substitution und semantische Folgerung]
Eine \emph{Substitution} ist ein Endomorphismus $\sigma:X^+\to X^+$;
jede Variable wird durch ein nichtleeres Wort ersetzt.
Die Abzählbarkeit von $X$ beschränkt weder die Mächtigkeit
noch die Erzeugerzahl der betrachteten Halbgruppen.
Für eine Menge $\Sigma$ von Wortidentitäten bedeutet
$\Sigma\models u=v$, daß jede Halbgruppe, die sämtliche
Identitäten aus $\Sigma$ erfüllt, auch $u=v$ erfüllt.
Diese Schreibweise bezeichnet die \emph{semantische Folgerung}.
\end{definition}

\begin{theorem}[Birkhoffs Vollständigkeitssatz für Wortidentitäten]
\label{b28guide:thm:equational-completeness}
Für jede Menge $\Sigma\subseteq X^+\times X^+$ und $u,v\in X^+$
gilt $\Sigma\models u=v$ genau dann, wenn es eine endliche Folge
$u=w_0,\ldots,w_n=v$ gibt, deren Schritte jeweils die Form
\[
 w_i=p\,\sigma(a)\,q,\qquad
 w_{i+1}=p\,\sigma(b)\,q
\]
mit $(a,b)\in\Sigma\cup\Sigma^{-1}$, einer Substitution $\sigma$
und Kontexten $p,q\in X^*$ haben.
Dabei ist $\Sigma^{-1}=\{(b,a):(a,b)\in\Sigma\}$ die umgekehrte
Relation; die Schreibweise bezeichnet keine multiplikativen Inversen.
Die Daten dürfen von Schritt zu Schritt wechseln; $n=0$ ist zugelassen.
Insbesondere folgt aus $\Sigma\models u=v$ bereits
$\Sigma_0\models u=v$ für eine endliche Teilmenge
$\Sigma_0\subseteq\Sigma$.

Der Satz charakterisiert die logische Folgerung aus beliebigen
Wortgleichungen. Die Endlichkeit betrifft allein eine Ableitung
und die darin verwendeten Gleichungen; weder $\Sigma$ noch die
betrachteten Halbgruppen müssen endlich sein.
Leere Kontexte lassen Faktoren entfallen; sie setzen keine
neutralen Elemente in den Halbgruppen voraus. \cite{burris}
\end{theorem}
